% GENERATED FILE. Edit canonical lessons, then run npm run generate:books.
% canonical-source-hash: fnv1a64:09067bf70182c6b9
% canonical-lessons: GE-C322-zutat, GE-C322-prise, GE-C322-essloeffel, GE-C322-abmessen, GE-C322-schuessel

\chapter{Die Zutaten}
\label{ch:ge-die-zutaten}
\hlchaptermodality{322}

\begin{chapteropening}
\textbf{By the end of this chapter:} \emph{I can get ready to cook: an ingredient, a pinch, a tablespoon, to measure out and a bowl.}

Complete the last lesson of chapter 322: I can get ready to cook: an ingredient, a pinch, a tablespoon, to measure out and a bowl.
\end{chapteropening}

\section[die Zutat]{die Zutat — the ingredient}

\label{lesson:GE-C322-zutat}



\begin{quote}
Before the new one: say the German for to come into being, then the German for to return.
\end{quote}

\subsection*{You'll want to know: die Zutat}
\textbf{die Zutat} — \textquotedblleft{}the ingredient\textquotedblright{}.

\begin{quote}
Für den Kuchen braucht man nur fünf Zutaten.
\end{quote}

\begin{quote}
\emph{For the cake you only need five ingredients.}
\end{quote}

Plural: \emph{die Zutaten}.

\begin{cousinweb}
\emph{Zutat} is built on \emph{zutun} \textquotedblleft{}to add\textquotedblright{}, from \emph{tun} \textquotedblleft{}to do, to put\textquotedblright{}, Old High German \emph{tuon} (EtymDB): what is put in.
\end{cousinweb}

\begin{grammarlens}[title={What you've built}]
One more word for telling, explaining and describing at B1.
\end{grammarlens}

\subsection*{Your turn}
\begin{itemize}
\raggedright
  \item \emph{Say it:} \emph{die Zutat}
  \item \emph{Say it:} \emph{die Zutat}, once more
  \item \emph{Say it:} a sentence of your own with \emph{die Zutat}
  \item \emph{Recall:} say the German for to come into being, then the German for to return, then say \emph{die Zutat} again
\end{itemize}

\begin{tcolorbox}[breakable,colback=teal!4,colframe=teal!35!black,title={Before you move on}]
What does die Zutat mean? (\textquotedblleft{}The ingredient\textquotedblright{}.) Say it once more.
\end{tcolorbox}
\section[die Prise]{die Prise — the pinch (of salt or spice)}

\label{lesson:GE-C322-prise}



\begin{quote}
Before the new one: say the German for to return, then the German for the ingredient.
\end{quote}

\subsection*{You'll want to know: die Prise}
\textbf{die Prise} — \textquotedblleft{}the pinch (of salt or spice)\textquotedblright{}.

\begin{quote}
Dann kommt eine Prise Salz dazu.
\end{quote}

\begin{quote}
\emph{Then a pinch of salt goes in.}
\end{quote}

Plural: \emph{die Prisen}. \emph{Eine Prise Zucker}.

\begin{cousinweb}
\emph{Prise} was borrowed from French \emph{prise} \textquotedblleft{}a taking, a hold\textquotedblright{} (EtymDB): as much as you take between two fingers.
\end{cousinweb}

\begin{grammarlens}[title={What you've built}]
One more word for telling, explaining and describing at B1.
\end{grammarlens}

\subsection*{Your turn}
\begin{itemize}
\raggedright
  \item \emph{Say it:} \emph{die Prise}
  \item \emph{Say it:} \emph{die Prise}, once more
  \item \emph{Say it:} a sentence of your own with \emph{die Prise}
  \item \emph{Recall:} say the German for to return, then the German for the ingredient, then say \emph{die Prise} again
\end{itemize}

\begin{tcolorbox}[breakable,colback=teal!4,colframe=teal!35!black,title={Before you move on}]
What does die Prise mean? (\textquotedblleft{}The pinch (of salt or spice)\textquotedblright{}.) Say it once more.
\end{tcolorbox}
\section[der Esslöffel]{der Esslöffel — the tablespoon}

\label{lesson:GE-C322-essloeffel}



\begin{quote}
Before the new one: say the German for the ingredient, then the German for the pinch.
\end{quote}

\subsection*{You'll want to know: der Esslöffel}
\textbf{der Esslöffel} — \textquotedblleft{}the tablespoon\textquotedblright{}.

\begin{quote}
Man nimmt zwei Esslöffel Zucker.
\end{quote}

\begin{quote}
\emph{You take two tablespoons of sugar.}
\end{quote}

Plural: \emph{die Esslöffel}. Short: \emph{EL}. A teaspoon is \emph{ein Teelöffel}.

\begin{cousinweb}
\emph{Esslöffel} joins \emph{essen} \textquotedblleft{}to eat\textquotedblright{} and \emph{Löffel} \textquotedblleft{}spoon\textquotedblright{}, from Proto-Germanic \emph{lapilaz} \textquotedblleft{}spoon\textquotedblright{} (EtymDB).
\end{cousinweb}

\begin{grammarlens}[title={What you've built}]
One more word for telling, explaining and describing at B1.
\end{grammarlens}

\subsection*{Your turn}
\begin{itemize}
\raggedright
  \item \emph{Say it:} \emph{der Esslöffel}
  \item \emph{Say it:} \emph{der Esslöffel}, once more
  \item \emph{Say it:} a sentence of your own with \emph{der Esslöffel}
  \item \emph{Recall:} say the German for the ingredient, then the German for the pinch, then say \emph{der Esslöffel} again
\end{itemize}

\begin{tcolorbox}[breakable,colback=teal!4,colframe=teal!35!black,title={Before you move on}]
What does der Esslöffel mean? (\textquotedblleft{}The tablespoon\textquotedblright{}.) Say it once more.
\end{tcolorbox}
\section[abmessen]{abmessen — to measure out}

\label{lesson:GE-C322-abmessen}



\begin{quote}
Before the new one: say the German for the pinch, then the German for the tablespoon.
\end{quote}

\subsection*{You'll want to know: abmessen}
\textbf{abmessen} — \textquotedblleft{}to measure out\textquotedblright{}.

\begin{quote}
Sie maß das Mehl genau ab.
\end{quote}

\begin{quote}
\emph{She measured out the flour exactly.}
\end{quote}

Separable, strong: \emph{misst ab}, \emph{maß ab}, \emph{hat abgemessen}.

\begin{cousinweb}
\emph{abmessen} is built on \emph{messen} \textquotedblleft{}to measure\textquotedblright{}, from Old High German \emph{mezzen}, Proto-Germanic \emph{metaną} (EtymDB), with \emph{ab}.
\end{cousinweb}

\begin{grammarlens}[title={What you've built}]
One more word for telling, explaining and describing at B1.
\end{grammarlens}

\subsection*{Your turn}
\begin{itemize}
\raggedright
  \item \emph{Say it:} \emph{abmessen}
  \item \emph{Say it:} \emph{abmessen}, once more
  \item \emph{Say it:} a sentence of your own with \emph{abmessen}
  \item \emph{Recall:} say the German for the pinch, then the German for the tablespoon, then say \emph{abmessen} again
\end{itemize}

\begin{tcolorbox}[breakable,colback=teal!4,colframe=teal!35!black,title={Before you move on}]
What does abmessen mean? (\textquotedblleft{}To measure out\textquotedblright{}.) Say it once more.
\end{tcolorbox}
\section[die Schüssel]{die Schüssel — the bowl}

\label{lesson:GE-C322-schuessel}



\begin{quote}
Before the new one: say the German for the tablespoon, then the German for to measure out.
\end{quote}

\subsection*{You'll want to know: die Schüssel}
\textbf{die Schüssel} — \textquotedblleft{}the bowl\textquotedblright{}.

\begin{quote}
Er gab die Eier in eine große Schüssel.
\end{quote}

\begin{quote}
\emph{He put the eggs into a big bowl.}
\end{quote}

Plural: \emph{die Schüsseln}.

\begin{cousinweb}
\emph{Schüssel} comes from Middle High German \emph{schüzzel} (EtymDB), usually explained as an early loan from Latin \emph{scutella} \textquotedblleft{}a small flat dish, a tray\textquotedblright{} (EtymDB).
\end{cousinweb}

\begin{grammarlens}[title={What you've built}]
One more word for telling, explaining and describing at B1.
\end{grammarlens}

\subsection*{Your turn}
\begin{itemize}
\raggedright
  \item \emph{Say it:} \emph{die Schüssel}
  \item \emph{Say it:} \emph{die Schüssel}, once more
  \item \emph{Say it:} a sentence of your own with \emph{die Schüssel}
  \item \emph{Recall:} say the German for the tablespoon, then the German for to measure out, then say \emph{die Schüssel} again
\end{itemize}

\begin{tcolorbox}[breakable,colback=teal!4,colframe=teal!35!black,title={Before you move on}]
What does die Schüssel mean? (\textquotedblleft{}The bowl\textquotedblright{}.) Say it once more.
\end{tcolorbox}
